\label{chap.2}

\section{Introduction}

Battery management for lithium-ion cells sits at the convergence of
electrochemical modelling, control theory, state estimation, and
machine learning.  This chapter surveys the four principal paradigms
in chronological and thematic order, identifies the key gaps that
motivate the present work, and maps each contribution of this thesis
to the gap it addresses.

\section{Classical Charging Protocols}

The industry-standard CC-CV protocol applies a constant current
until the terminal voltage reaches the cut-off threshold, then holds
the voltage constant while the current decays to a minimum cut-off
\cite{bloom2001accelerated}.  Variants such as multi-stage constant
current (MCC) and pulse charging have been proposed to reduce
polarisation losses and thermal excursions
\cite{purushothaman2006multistage}.

The fundamental limitation of all classical protocols is their
open-loop, parameter-agnostic nature.  They do not adapt to
cell-to-cell variability in internal resistance, ambient temperature,
or the degradation state of the cell.  As the series resistance
$R_0$ grows with ageing, a fixed-current protocol produces
increasingly large voltage drops, causing the CV transition to occur
at progressively lower SOC values, effectively truncating each charge
cycle.

\section{Model Predictive Control for Battery Charging}

MPC-based fast charging formulates a finite-horizon constrained
optimisation problem over an electrochemical or equivalent circuit
model of the battery \cite{torchio2016lionsimba}.  The Single
Particle Model (SPM) or Doyle-Fuller-Newman (DFN/P2D) model provides
physics-based state equations; the MPC minimises charging time
subject to internal potential and temperature constraints.

Torchio et al.~\cite{torchio2016lionsimba} demonstrated that
NMPC with a P2D model can reduce charging time by 20--30\% relative
to CC-CV while respecting lithium plating and temperature limits.
However, two critical limitations apply: (1) P2D numerical
integration carries prohibitive real-time computational cost, and
(2) the model parameters ($R_0$, diffusion coefficients, heat
generation) are assumed fixed.  As the cell ages, the growing
$R_0$ causes the MPC to underestimate voltage and overshoot
current, causing unmodelled plating at high SOC.

\section{Single-Agent Deep Reinforcement Learning}

Deep Reinforcement Learning (DRL) has been applied to battery fast
charging as a way to learn current policies directly from
simulator interactions, sidestepping the need for explicit
electrochemical models \cite{park2020fastcharging, wu2021dqn}.
Park et al.~\cite{park2020fastcharging} used a Deep Deterministic
Policy Gradient (DDPG) agent to discover optimal multi-step
current profiles in an equivalent circuit model, reporting an
improvement in charging time over MCC protocols.

The primary weakness of single-agent DRL is the \emph{objective
entanglement} problem.  The scalar reward
$R_t = -(w_t \cdot \Delta t + w_d \cdot D_t + w_e \cdot E_t + w_v
\cdot V_{\mathrm{pen},t})$ aggregates charging time, degradation,
energy cost, and voltage violations with fixed weights.  The resulting
policy is acutely sensitive to these weight ratios: small
perturbations cause the learned policy to either sacrifice cell
health or miss charging deadlines entirely.  Wu et al.~\cite{wu2021dqn}
introduced a DQN with heuristic thermal throttling, adding a rule-based
override on top of the neural policy.  While this reduces thermal
violations, it introduces chattering instability near the threshold.

\section{Multi-Agent Reinforcement Learning in Battery Systems}

Multi-Agent Reinforcement Learning (MARL) has been applied
predominantly to multi-cell balancing in battery packs and energy
market trading for large-scale battery energy storage systems (BESS)
\cite{cao2020marl_energy, li2022marl_grid, zhang2022maddpg}.  In
these formulations, each spatial entity (cell, inverter, generator)
is assigned a dedicated agent, and agents learn cooperative or
competitive policies.

The present work differs fundamentally: rather than assigning agents
to spatial entities, we assign agents to \emph{functional control
objectives} of a single pack.  This \emph{functional decomposition}
paradigm preserves the modularity advantages of multi-agent systems
while eliminating inter-cell coupling complexity.

\section{Safe Reinforcement Learning}

Safe RL addresses the critical failure mode of neural policies: their
inability to \emph{guarantee} constraint satisfaction.  Achiam et
al.~\cite{achiam2017cpo} proposed Constrained Policy Optimisation
(CPO), adding a Lagrangian constraint on cumulative cost.  Zou et
al.~\cite{zou2023cbf} combined DDPG with a Control Barrier Function
(CBF) safety filter, projecting unsafe actions onto the safe manifold
via a real-time quadratic program.

Both approaches share a fundamental vulnerability: the safety
guarantee depends on accurate state observation.  When a temperature
sensor is stuck or biased low by $-5\,^\circ\mathrm{C}$, the observed
state appears safe while the true cell is overheating.  The CBF
constraint is then satisfied on a false observation, permitting
thermal runaway.  Our deterministic analytical supervisor addresses
this by incorporating $R_0$ margins and a thermal plausibility
residual monitor operating on model predictions, not sensor readings
alone.

\section{State and Parameter Estimation Under Ageing}

Accurate state-of-charge (SOC) and state-of-health (SOH) estimation
is essential for all BMS paradigms.  The Extended Kalman Filter (EKF)
is the canonical approach \cite{plett2004ekf}, with extensions to dual
EKF (simultaneous estimation of SOC and $R_0$) and sigma-point
variants \cite{hu2012ekf}.  Wang et al.~\cite{wang2017dual_ekf}
proposed a joint EKF for SOC--SOH estimation, improving accuracy under
degradation.

The key gap in existing literature is \emph{initialisation under
aged conditions}: estimators that initialise $\rho_0 = 1.0$ (fresh
cell assumption) for an aged cell with $R_0 \approx 1.5\times$ the
nominal value will systematically under-predict terminal voltage,
causing premature current derating and missed deadlines.  Our 4-state
EKF initialises $\rho_0 = r_0\text{\_scale}(\widehat{\mathrm{SOH}})$,
which eliminates this bias.

\section{Empirical Grounding and RUL Evaluation}

Rigorous experimental validation is the weakest link in battery
control literature.  Most DRL charging papers use synthetic ECM
simulators fitted to fresh cells, without auditing the raw
dataset for logging artefacts or charge-phase telemetry
completeness.  Studies using the NASA PCoE dataset
\cite{saha2008battery, fricke2023nasa} frequently misinterpret
the logging definitions, misclassify relay-open transient collapses
as internal short circuits, and assume all 26 packs experienced
classical capacity-fade EOL.

RUL evaluation is equally fragile: treating RUL as a standard
regression target over a small sample of fewer than 15 failures
without accounting for right-censored packs leads to severely
optimistic error estimates \cite{saxena2015remaining}.

\section{Comparative Synthesis of Related Works}

To establish a clear structural overview of where existing approaches fall
short, Table~\ref{tab:lit_comparison} evaluates representative battery
charging, management, and prognostics literature across six foundational
axes: control paradigm, multi-objective formulation, aging awareness,
safety guarantees, fault robustness, and empirical validation grounding.

\begin{table}[htbp]
  \centering
  \caption{Systematic comparison of related battery management literature against the proposed MA-BMS framework.}
  \label{tab:lit_comparison}
  \resizebox{\columnwidth}{!}{%
  \begin{tabular}{lp{2.8cm}p{3.2cm}cccp{3.5cm}}
    \toprule
    \textbf{Study / Reference} & \textbf{Control Paradigm} & \textbf{Objective Formulation} &
    \textbf{Aging ($R_0$)} & \textbf{Hard Safety} & \textbf{Fault Robust} &
    \textbf{Empirical Grounding} \\
    \midrule
    Bloom et al.~\cite{bloom2001accelerated} &
      Classical CC-CV & Single objective (time) & Static & No (heuristic) & No &
      Laboratory cycler \\
    Torchio et al.~\cite{torchio2016lionsimba} &
      Electrochemical NMPC & Constrained min-time & Static (fresh) & Partial (model-brittle) & No &
      Synthetic P2D simulation \\
    Park et al.~\cite{park2020fastcharging} &
      Single-Agent DDPG & Entangled scalar reward & Static fixed & No (soft penalty) & No &
      Synthetic ECM \\
    Wu et al.~\cite{wu2021dqn} &
      DQN + Heuristic Rule & Scalar reward + override & Static fixed & Partial (rule derating) & No &
      Synthetic ECM \\
    Zou et al.~\cite{zou2023cbf} &
      Safe DRL via CBF & Scalar reward + QP filter & Static fresh & Partial (bias-vulnerable) & No &
      Simulation only \\
    Zhang et al.~\cite{zhang2022maddpg} &
      Spatial MADDPG & Spatial cell balancing & Static & No (soft penalty) & No &
      Synthetic multi-cell \\
    Wang et al.~\cite{wang2017dual_ekf} &
      Joint Dual EKF & State estimation only & Dynamic tracking & N/A (observer only) & Partial &
      Pulsed cell tests \\
    \midrule
    \textbf{Proposed MA-BMS} &
      \textbf{Functional Multi-Agent} & \textbf{Decoupled Proposers} &
      \textbf{Adaptive $R_0(\mathrm{SOH})$} & \textbf{Yes (Analytical $I_{\mathrm{safe}}$)} &
      \textbf{Yes (Monitor)} & \textbf{NASA ALT (4.9M rows, validated} \\
    \textbf{(This Thesis)} &
      \textbf{+ Safety Supervisor} & \textbf{+ Dirichlet Tuning} &
      \textbf{$k_{R_0} \sim \mathcal{U}(1.5, 4.5)$} & \textbf{0\% hard violations} &
      \textbf{7 fault types} & \textbf{thermal cooling \& OCV)} \\
    \bottomrule
  \end{tabular}%
  }
\end{table}

\section{Summary of Gaps and Thesis Positioning}

Table~\ref{tab:lit_gaps} summarises the key gaps and their
correspondence to contributions in this thesis.

\begin{longtable}{>{\raggedright}p{4cm}
                  >{\raggedright}p{4.5cm}
                  >{\raggedright\arraybackslash}p{4cm}}
\caption{Summary of literature gaps and corresponding contributions.}
\label{tab:lit_gaps}\\
\toprule
\textbf{Gap} & \textbf{Limitation in Literature} &
\textbf{Contribution in This Thesis} \\
\midrule
\endfirsthead
\multicolumn{3}{c}{\tablename\ \thetable\ -- \textit{Continued}}\\
\toprule
\textbf{Gap} & \textbf{Limitation in Literature} &
\textbf{Contribution in This Thesis} \\
\midrule
\endhead
\bottomrule
\endfoot
Static $R_0$ in controllers &
MPC and DRL controllers calibrated on fresh cells fail under aged
resistance ($1.57\times$ growth) &
Aging-dependent $R_0(\mathrm{SOH})$ in simulator; adaptive EKF
initialisation \\[4pt]
Objective entanglement &
Single scalar reward collapses under weight perturbation &
Functional agent decomposition + Dirichlet weight sensitivity analysis
\\[4pt]
No deterministic safety guarantee &
Neural policies violate hard limits under sensor faults &
Analytical safety supervisor: 0 violations across 7 fault types \\[4pt]
Dataset audit gaps &
Misclassified EOL events; unverified charge telemetry &
Full mode audit; CC-only charging confirmed; corrected EOL labels \\[4pt]
Na\"{\i}ve RUL evaluation &
Standard regression ignores censoring and small sample &
Interval-censored Kaplan-Meier + LOPO Ridge regression \\
\end{longtable}
